\section{第3章 微分中值定理与导数的应用}
\subsection{中值定理}
\begin{exercise}{289}{罗尔定理的条件与结论}
分别判断函数 $f(x)=|x|$、$f(x)=x^2-1$、$f(x)=x\mathrm e^{-x}$、$f(x)=1/x$ 在区间 $[-1,1]$、$[-1,1]$、$[0,1]$、$[-1,1]$ 上是否满足罗尔定理的全部条件，逐项指出依据。对满足条件的函数求全部罗尔点；对不满足条件的函数，判断开区间内是否仍有导数为零的点。
\end{exercise}
\begin{exercise}{290}{罗尔点的存在性与唯一性}
验证 $f(x)=\ln\sin x$ 在 $[\pi/6,5\pi/6]$ 上满足罗尔定理，求出全部罗尔点。若将右端点改为 $b\in(\pi/6,\pi)$，求使两个端点函数值相等的全部 $b$，并说明相应罗尔点是否唯一。
\end{exercise}
\begin{exercise}{291}{辅助函数与导数方程}
设 $f\in C[0,1]$，在 $(0,1)$ 可导，$f(0)=0$、$f(1)=\pi/4$。证明存在 $\xi\in(0,1)$ 使 $(1+\xi^2)f'(\xi)=1$，写出所用辅助函数。将端点条件改为 $f(1)=c$ 时，求使同一辅助函数的两个端点值相等的 $c$，并解释该方法为何依赖此值。
\end{exercise}
\begin{exercise}{292}{多项式零点与原函数}
设实数 $a_0,\ldots,a_n$ 满足 $\sum_{k=0}^n a_k/(k+1)=0$。不用定积分，证明多项式 $P(x)=\sum_{k=0}^n a_kx^k$ 在 $(0,1)$ 内有零点。再说明若 $P$ 不恒为零，则它在 $(0,1)$ 内必同时取正值和负值。
\end{exercise}
\begin{exercise}{293}{极限条件与二阶罗尔定理}
设 $f$ 在包含 $[0,1]$ 的开区间内二阶可导，且 $\lim_{x\to0}f(x)/x=0$、$f(1)=0$。由极限条件确定 $f(0)$ 与 $f'(0)$，再证明存在 $\xi\in(0,1)$ 使 $f''(\xi)=0$。若仅保留 $f(0)=f(1)=0$，举例说明二阶导数为零的结论是否仍成立。
\end{exercise}
\begin{exercise}{294}{端点高阶零点与连续使用罗尔定理}
设 $f$ 在包含 $[0,1]$ 的开区间内三阶可导，$f(0)=f(1)=0$，令 $F(x)=x^3f(x)$。求 $F(0)$、$F'(0)$，证明在 $(0,1)$ 内至少存在一点使 $F''$ 为零。进一步证明 $F'''$ 在 $(0,1)$ 内也至少有一个零点，明确每次应用罗尔定理的区间。
\end{exercise}
\begin{exercise}{295}{端点匹配与指数辅助函数}
设 $f\in C[0,1]$，在 $(0,1)$ 可导，$f(0)=1$、$f(1)=\mathrm e^{-1}$。证明存在 $\xi\in(0,1)$ 使 $f'(\xi)+\mathrm e^{-\xi}=0$。若另知 $f''(x)>\mathrm e^{-x}$ 对所有 $x\in(0,1)$ 成立，证明该点唯一；此问假设 $f$ 在 $(0,1)$ 二阶可导。
\end{exercise}
\begin{exercise}{296}{乘积辅助函数与参数推广}
设 $a>0$，$f\in C[0,a]$，在 $(0,a)$ 可导，且 $f(a)=0$。证明存在 $\xi\in(0,a)$ 使 $f(\xi)+\xi f'(\xi)=0$。对任意给定 $p>0$，进一步证明存在 $\eta\in(0,a)$ 使 $pf(\eta)+\eta f'(\eta)=0$，并说明为何无需假设 $x^pf(x)$ 在 $0$ 处可导。
\end{exercise}
\begin{exercise}{297}{加权导数零点}
设 $f\in C[0,1]$，在 $(0,1)$ 可导，且 $f(1)=0$。证明存在 $\xi\in(0,1)$ 使 $f'(\xi)=-2f(\xi)/\xi$。讨论若把区间改为 $[a,1]$，其中 $0<a<1$，且仍只给出 $f(1)=0$，结论是否必然成立；给出证明或反例。
\end{exercise}
\begin{exercise}{298}{对称乘积与对数导数}
设 $f\in C[0,1]$，在 $(0,1)$ 可导，$f(0)=0$，且 $f(x)\ne0$ 对 $0<x<1$ 成立。证明存在 $\xi\in(0,1)$ 使
\[\frac{f'(\xi)}{f(\xi)}=\frac{f'(1-\xi)}{f(1-\xi)}.\]
分别用对称乘积构造法与取 $\xi=1/2$ 的方法论证，并说明上述条件中哪些对后一个论证并非必要。
\end{exercise}
\begin{exercise}{299}{平方辅助函数与非零分母}
设 $f\in C[0,1]$，在 $(0,1)$ 可导，$f(0)=f(1/2)=1$、$f(1)=2$。证明存在 $\xi\in(0,1)$ 使 $2f(\xi)f'(\xi)=\mathrm e^{\xi}$，从而 $f'(\xi)=\mathrm e^{\xi}/[2f(\xi)]$。证明中须说明辅助函数为何取得相同的两个函数值，以及所得点上 $f(\xi)\ne0$。
\end{exercise}
\begin{exercise}{300}{离散函数值约束与驻点}
设 $f\in C[0,3]$，在 $(0,3)$ 可导，且 $f(0)+f(1)+f(2)=3$、$f(3)=1$。证明存在 $\xi\in(0,3)$ 使 $f'(\xi)=0$。再判断能否保证这样的点位于 $(0,2)$，并用具体函数支持结论。
\end{exercise}
\begin{exercise}{301}{介值定理与指数辅助函数}
设 $f\in C[0,1]$，在 $(0,1)$ 可导，$f(0)=f(1)=0$、$f(1/2)=1$。证明存在 $\eta\in(1/2,1)$ 使 $f(\eta)=\eta$。对任意实数 $\lambda$，证明存在 $\xi\in(0,\eta)$ 使 $f'(\xi)-\lambda f(\xi)+\lambda\xi=1$，并写出与所选 $\eta$ 对应的辅助函数。
\end{exercise}
\begin{exercise}{302}{三零点与二阶导数}
设抛物线 $p(x)=-x^2+Bx+C$ 的两个零点为 $a<b$，$f$ 在包含 $[a,b]$ 的开区间内二阶可导，且 $f(a)=f(b)=0$。若存在 $c\in(a,b)$ 使 $f(c)=p(c)$，证明存在 $\xi\in(a,b)$ 使 $f''(\xi)+2=0$。说明若交点仅允许为端点，为何结论不再得到保证。
\end{exercise}
\begin{exercise}{303}{拉格朗日中值定理的适用条件}
分别判断 $\ln(1+x^2)$、$|x|$、$1/\sqrt{1-x^2}$、$\sin(1/x)$ 在 $[-1,1]$ 上是否满足拉格朗日中值定理的条件，指出每项缺失或满足的条件。对第一项求所有中值点；对第二项直接判断中值等式是否存在解。
\end{exercise}
\begin{exercise}{304}{端点连续性不可省略}
设 $a<b$，$f$ 在 $[a,b]$ 有定义且在 $(a,b)$ 可导。判断下列四个断言是否必然成立，错误断言须给出反例：若 $f(a)f(b)<0$，则内部有驻点；若 $f(a)=f(b)$，则内部有驻点；存在 $\xi\in(a,b)$ 使 $f(b)-f(a)=f'(\xi)(b-a)$；对任意 $\xi\in(a,b)$ 都有 $\lim_{x\to b^-}[f(x)-f(\xi)]=0$。说明加上 $f\in C[a,b]$ 后，哪些断言变为真。
\end{exercise}
\begin{exercise}{305}{中值定理证明三角不等式}
用拉格朗日中值定理证明 $\sin x/x>\cos x$ 对 $0<x<\pi$ 成立。再判断该不等式在 $-\pi<x<0$ 是否成立，说明奇偶性在推广中的作用。
\end{exercise}
\begin{exercise}{306}{对数增量的双侧估计}
设 $a>b>0$。用中值定理证明
\[\frac{a-b}{a}<\ln\frac ab<\frac{a-b}{b}.\]
固定 $b$，求 $a\to b^+$ 时 $\ln(a/b)$ 与这两个界的比值，说明严格不等号为何不妨碍一阶近似。
\end{exercise}
\begin{exercise}{307}{对数不等式与余项控制}
用中值定理证明 $x/(1+x)<\ln(1+x)<x$ 对 $x>0$ 成立。利用这一结论证明 $0<x-\ln(1+x)<x^2/(1+x)$，并判断所得上界能否直接用于确定 $[x-\ln(1+x)]/x^2$ 的精确极限。
\end{exercise}
\begin{exercise}{308}{指数函数的切线比较}
用中值定理证明 $\mathrm e^x>\mathrm e x$ 对 $x>1$ 成立。进一步确定全部实数 $x$ 上 $\mathrm e^x-\mathrm e x$ 的最小值及等号条件，并解释两部分论证的联系。
\end{exercise}
\begin{exercise}{309}{导数恒等式与端点延拓}
在 $(-1,1)$ 上用求导与中值定理证明 $\arcsin x+\arccos x=\pi/2$，再将恒等式延拓至 $[-1,1]$。说明在端点不能直接使用通常的反三角函数求导公式，仍可得到等式的依据。
\end{exercise}
\begin{exercise}{310}{反三角函数的分段恒等式}
证明 $2\arctan x+\arcsin[2x/(1+x^2)]=\pi$ 对 $x\ge1$ 成立。进一步求该表达式在 $0\le x\le1$ 上的简化式，说明反正弦主值范围为何导致分段。
\end{exercise}
\begin{exercise}{311}{乘积函数的中值等式}
设 $f\in C[a,b]$，在 $(a,b)$ 可导，$a<b$。证明存在 $\xi\in(a,b)$ 使
\[\frac{bf(b)-af(a)}{b-a}=f(\xi)+\xi f'(\xi).\]
取 $f(x)=x^2$、$[a,b]=[1,2]$ 求全部满足等式的点，并检查它们确实位于开区间内。
\end{exercise}
\begin{exercise}{312}{导数单调与次可加性}
设 $c>0$，$f\in C[0,c]$，在 $(0,c)$ 可导，$f'$ 单调不增，$f(0)=0$。对 $0\le a\le b\le a+b\le c$，用拉格朗日中值定理证明 $f(a+b)\le f(a)+f(b)$。若 $f'$ 严格递减，讨论 $a=0$ 与 $a>0$ 两种情况下等号能否成立。
\end{exercise}
\begin{exercise}{313}{柯西中值定理与导数非零条件}
对 $f(x)=\sin x$、$g(x)=1+\cos x$ 在 $[0,\pi/2]$ 上验证柯西中值定理，求全部满足
\[\frac{f(\pi/2)-f(0)}{g(\pi/2)-g(0)}=\frac{f'(\xi)}{g'(\xi)}\]
的 $\xi$。说明 $g'(0)=0$ 是否妨碍应用定理，以及分母为何在所求点处非零。
\end{exercise}
\begin{exercise}{314}{倒数变量与柯西中值定理}
设 $0<a<b$，$f\in C[a,b]$ 且在 $(a,b)$ 可导。用柯西中值定理证明存在 $\xi\in(a,b)$ 使
\[\frac{af(b)-bf(a)}{a-b}=f(\xi)-\xi f'(\xi).\]
明确选取的两个辅助函数，并指出 $a>0$ 在此证明中保证了什么。
\end{exercise}
\begin{exercise}{315}{幂函数参照与加权中值点}
设 $0<a<b$，$f\in C[a,b]$ 且在 $(a,b)$ 可导。证明存在 $\xi\in(a,b)$ 使
\[\frac{f(b)-f(a)}{b-a}=(a^2+ab+b^2)\frac{f'(\xi)}{3\xi^2}.\]
将区间改为 $[-1,1]$，判断上述含商等式是否仍对所有可导函数成立，并说明原证明中哪一条件失效。
\end{exercise}
\begin{exercise}{316}{对数复合函数的中值点}
用拉格朗日中值定理证明存在 $\xi\in(1,\mathrm e)$ 使 $\sin1=\cos(\ln\xi)$，写出中值定理所作用的区间与函数。利用余弦在 $(0,1)$ 上的单调性，证明这样的 $\xi$ 唯一。
\end{exercise}
\begin{exercise}{317}{带参数的中值恒等式}
设 $f\in C[0,1]$，在 $(0,1)$ 可导，$f(0)=0$、$f(1)=1$。证明对每个实数 $\lambda$ 都存在 $\xi\in(0,1)$ 使 $f'(\xi)+\lambda[f(\xi)-\xi]=1$。若要求同一个 $\xi$ 对所有 $\lambda$ 同时适用，现有条件是否足够？给出证明或反例。
\end{exercise}
\begin{exercise}{318}{线性插值误差与二阶中值定理}
设 $f\in C^2[a,b]$，$a<b$，$L(x)=[(b-x)f(a)+(x-a)f(b)]/(b-a)$。证明对每个 $x\in(a,b)$ 存在 $\xi\in(a,b)$ 使 $f(x)-L(x)=\tfrac12f''(\xi)(x-a)(x-b)$。由此在 $|f''|\le M$ 的条件下给出整个 $[a,b]$ 上的统一误差界，并说明界中常数能否改小。
\end{exercise}
\subsection{洛必达法则}
\begin{exercise}{319}{高阶分母与单侧无穷极限}
讨论 $\lim_{x\to1}(x^3-3x^2+2)/(x^3-x^2-x+1)$ 是否存在，并分别求左右极限。先判断不定式类型，再用因式分解或洛必达法则分析零点阶数，说明不能在分母求导后为零时直接把该步写成有限数值。
\end{exercise}
\begin{exercise}{320}{基本对数极限与二阶修正}
用洛必达法则求 $\lim_{x\to1}\ln x/(x-1)$，核对使用条件。再求 $\lim_{x\to1}[\ln x-(x-1)]/(x-1)^2$，说明再次求导前为何必须重新检查是否仍为不定式。
\end{exercise}
\begin{exercise}{321}{混合函数的一阶极限}
用洛必达法则求 $\lim_{x\to1}[x^3-1+\ln x]/(\mathrm e^x-\mathrm e)$。若将分子中的 $\ln x$ 改为 $\alpha\ln x$，求使极限为零的实参数 $\alpha$，再求该参数下整个商与 $x-1$ 的比值极限。
\end{exercise}
\begin{exercise}{322}{指数与多项式的增长比较}
求 $\lim_{x\to+\infty}(x^4+\mathrm e^x)/(x\mathrm e^x)$。进一步对实数 $p$ 讨论 $\lim_{x\to+\infty}(x^4+\mathrm e^x)/(x^p\mathrm e^x)$，用洛必达法则或等价增长比较给出分类，不能只凭“指数增长快”省略关键极限。
\end{exercise}
\begin{exercise}{323}{三阶零点与连续洛必达}
求 $\lim_{x\to0}(x-\sin x)/(x^2\sin x)$，每次使用洛必达法则前确认不定式。再求使 $x-\sin x-cx^3=o(x^3)$ 的 $c$，说明结果与所求极限之间的关系。
\end{exercise}
\begin{exercise}{324}{对数与任意正幂的比较}
设 $a>0$，求 $\lim_{x\to+\infty}\ln(x\ln x)/x^a$，说明此处 $x$ 应取在什么范围。进一步讨论 $a=0$、$a<0$ 时的极限，展示正参数限制对结论的作用。
\end{exercise}
\begin{exercise}{325}{根式零点与指数三角相消}
求 $\lim_{x\to0}(\mathrm e^x-\sin x-1)/(1-\sqrt{1-x^2})$，写出实定义邻域。再把分子中的 $\sin x$ 改为 $\alpha\sin x$，求使双侧有限极限存在的全部实数 $\alpha$，并求相应极限。
\end{exercise}
\begin{exercise}{326}{反正切差与对数三阶零点}
求 $\lim_{x\to0}(\arctan x-x)/\ln(1+2x^3)$，给出对数在零点附近有意义的依据。将分母改为 $\ln(1+2x^2)$ 后再求极限，并从零点阶数说明两者不同的原因。
\end{exercise}
\begin{exercise}{327}{多个无穷小因子的主阶}
求
\[\lim_{x\to0}\frac{\sin^2x-x^2\cos^2x}{\sin x(\mathrm e^{2x}-1)\ln(1+\tan^2x)}.\]
先分别确定分子与分母的最低非零阶，再选用洛必达法则或等价无穷小完成计算，说明为何分子中的两项不能直接分别替换为 $x^2$ 后相减。
\end{exercise}
\begin{exercise}{328}{参数导致的零点阶数变化}
设 $a^2+b^2\ne0$。求使
\[\lim_{x\to0}\frac{a\tan x+b(1-\cos x)}{c\ln(1-2x)+d(1-\mathrm e^{-x^2})}=2\]
成立的全部实参数 $(a,b,c,d)$，分别处理 $c\ne0$、$c=0,d\ne0$、$c=d=0$。据此判断 $b=4d$、$b=-4d$、$a=4c$、$a=-4c$ 中哪些是所有可行参数必满足的关系。
\end{exercise}
\begin{exercise}{329}{对数三角复合的二阶相消}
求 $\lim_{x\to0}(1-\cos x)[x-\ln(1+\tan x)]/\sin^4x$。先求 $[x-\ln(1+\tan x)]/x^2$ 的极限，再组合结果，检查对数真数在零点邻域的正性。
\end{exercise}
\begin{exercise}{330}{嵌套正弦的等价无穷小}
求 $\lim_{x\to0}[\sin x-\sin(\sin x)]\sin x/x^4$。明确内变量 $u=\sin x$ 的使用，并进一步求 $[\sin x-\sin(\sin x)]/x^3$ 的极限，说明这两个极限的联系。
\end{exercise}
\begin{exercise}{331}{无穷量相减的恒等化简}
求 $\lim_{x\to0}x^{-1}(1/\sin x-1/\tan x)$。先合并为一个分式，再用半角公式或洛必达法则求值；说明两个倒数项分别无界时为何不能单独取极限再作减法。
\end{exercise}
\begin{exercise}{332}{多项式与指数衰减的参数条件}
设 $n$ 为正整数、$\lambda>0$，求 $\lim_{x\to+\infty}x^n\mathrm e^{-\lambda x}$，给出能处理任意 $n$ 的论证。再讨论 $\lambda=0$ 与 $\lambda<0$ 时的极限，说明正参数假设的必要性。
\end{exercise}
\begin{exercise}{333}{零乘无穷与线性变量替换}
求 $\lim_{x\to1}(1-x)\tan(\pi x/2)$，分别说明左右两侧的乘积符号。通过令 $t=\pi(1-x)/2$ 化为基本三角极限，再用洛必达法则核验。
\end{exercise}
\begin{exercise}{334}{三角无穷量之差与可去间断}
求 $\lim_{x\to\pi/2}(\sec x-\tan x)$。先用共轭恒等式消去无穷量相减，给出左右极限，并判断原表达式在 $\pi/2$ 补值后是否可导；可导时求导数。
\end{exercise}
\begin{exercise}{335}{倒数差的补值与导数}
设 $h(x)=1/x-1/(\mathrm e^x-1)$，$x\ne0$。求 $\lim_{x\to0}h(x)$，补定义 $h(0)$ 使其连续，再判断补值后的函数在 $0$ 是否可导并求导数。每次合并分式后说明所用极限方法的条件。
\end{exercise}
\begin{exercise}{336}{二阶主部消去与偶性}
求 $\lim_{x\to0}[1/x^2-1/(x\tan x)]$。给出连续补值后函数在 $0$ 的导数，并说明偶性只有在可导性已经建立后才能用于确定导数值。
\end{exercise}
\begin{exercise}{337}{零的零次幂型极限}
求 $\lim_{x\to0^+}x^x$，通过取对数把它化为可应用洛必达法则的形式。进一步求 $\lim_{x\to0^+}(x^x-1)/(x\ln x)$，说明原底数必须取正的原因。
\end{exercise}
\begin{exercise}{338}{一的无穷次幂型与符号}
求 $\lim_{x\to1}x^{1/(1-x)}$，明确取对数后分母的符号。对实参数 $a$ 求 $\lim_{x\to1}x^{a/(1-x)}$，说明 $a=0$ 时与一般公式是否相容。
\end{exercise}
\begin{exercise}{339}{三角分母指数的幂极限}
求 $\lim_{x\to0}(1+3x)^{2/\sin x}$，给出底数为正的邻域。再求对实参数 $a,b$ 的 $\lim_{x\to0}(1+ax)^{b/\sin x}$，说明双侧取极限的合法性与指数为零的情形。
\end{exercise}
\begin{exercise}{340}{二阶对数之比的幂极限}
求 $\lim_{x\to0}(\cos x)^{1/\ln(1+x^2)}$。先确定取实数幂时底数为正的邻域，再处理对数之比。将分母改为 $\ln(1+ax^2)$、$a>0$ 后给出参数结果。
\end{exercise}
\begin{exercise}{341}{趋零底数与趋零指数}
求 $\lim_{x\to+\infty}(x^{1/x}-1)^{1/\ln x}$，先证明底数在所考虑范围为正。取对数后使用 $\mathrm e^u-1\sim u$，分开处理 $\ln x$ 与 $\ln\ln x$ 的增长，不能把指数趋于零直接解释为极限为 $1$。
\end{exercise}
\begin{exercise}{342}{双指数平均的幂极限}
设 $a,b>0$，求 $\lim_{x\to0}[(a^x+b^x)/2]^{3/x}$。进一步讨论一般实参数 $p$ 下指数改为 $p/x$ 的结果，并核验 $a=b$ 时可由原表达式直接得到同一结论。
\end{exercise}
\begin{exercise}{343}{复合指数差的无穷小阶数}
求实数 $n>0$，使 $x\to0^+$ 时 $\mathrm e^{\tan x}-\mathrm e^x$ 与 $x^n$ 为同阶无穷小，并求两者比值的非零极限。说明若把“同阶”加强为“等价”，是否仍能直接使用 $x^n$，应如何调整主项系数。
\end{exercise}
\begin{exercise}{344}{分段函数的连续补值与高阶信息}
设 $f(x)=[\sin2x+\mathrm e^{2ax}-1]/x$（$x\ne0$）、$f(0)=a$。求使 $f$ 在实数轴连续的全部实参数 $a$，再判断相应函数在 $0$ 是否可导并求导数。注意参数同时出现于非零分支和补定义值中，必须解出由连续性产生的参数方程。
\end{exercise}
\begin{exercise}{345}{幂函数的连续延拓}
在 $0<|x|<\pi/2$ 设 $f(x)=(\cos x)^{x^2}$，令 $f(0)=a$。求使其在 $0$ 连续的 $a$，再求 $\lim_{x\to0}[f(x)-a]/x^4$。比较将指数改为 $1/x^2$ 时的连续补值，说明指数中幂次对结果的影响。
\end{exercise}
\begin{exercise}{346}{端点奇异项的抵消}
设 $f(x)=1/(\pi x)+1/\sin\pi x-1/[\pi(1-x)]$，$1/2\le x<1$。求连续延拓值 $f(1)$，再求 $\lim_{x\to1^-}[f(x)-f(1)]/(x-1)$。令 $t=1-x$ 后明确哪些奇异项相消，不能分别代入取值。
\end{exercise}
\begin{exercise}{347}{迭代数列与相邻项的幂极限}
设 $0<x_1<\pi$，$x_{n+1}=\sin x_n$。证明数列从某项起严格递减且为正，求其极限；再分别求 $\lim_{n\to\infty}(x_{n+1}/x_n)^{1/x_n}$ 与 $\lim_{n\to\infty}(x_{n+1}/x_n)^{1/x_n^2}$，并说明即使 $x_1>\pi/2$ 也不影响论证。
\end{exercise}
\begin{exercise}{348}{洛必达法则失效时的原极限}
设 $f(x)=x+\sin x$、$g(x)=x$，$x>0$。求 $\lim_{x\to+\infty}f(x)/g(x)$，并证明 $f'(x)/g'(x)$ 没有极限。由此说明洛必达法则的导数比极限条件是应用定理的条件，能否把它当成原商有极限的必要条件。
\end{exercise}
\begin{exercise}{349}{参数主项消去与幂极限}
设 $a\in\mathbb R$，求 $\lim_{x\to0}(1+x+ax^2)^{1/x}$，再讨论 $\lim_{x\to0}(1+x+ax^2)^{1/x^2}\mathrm e^{-1/x}$。说明底数在零点附近为正，并用对数变换解释两问需要的不同展开阶数。
\end{exercise}
\subsection{泰勒公式}
\begin{exercise}{350}{两类泰勒余项与误差估计}
对任意正整数 $n$，写出 $\mathrm e^x$ 的 $n$ 阶麦克劳林公式，分别使用拉格朗日型余项与佩亚诺型余项。用三次麦克劳林多项式近似 $\mathrm e^{0.1}$，给出保证成立的绝对误差上界；说明仅写 $o(x^3)$ 能否直接提供这一数值误差界。
\end{exercise}
\begin{exercise}{351}{展开阶数与商的极限}
用泰勒公式求 $\lim_{x\to0}x\ln(1+x)/(\mathrm e^x-x-1)$，并求
\[\lim_{x\to0}\frac1x\left[\frac{x\ln(1+x)}{\mathrm e^x-x-1}-2\right].\]
说明求第二个极限时分子、分母分别至少需要展开到哪一阶。
\end{exercise}
\begin{exercise}{352}{三角函数的高阶消去}
用带佩亚诺余项的展开求 $\lim_{x\to0}(\sin x-x\cos x)/\tan^3x$。进一步求分子减去其三次主项后的五次主项，并说明仅将 $\sin x$ 换为 $x$ 为何会丢失所求信息。
\end{exercise}
\begin{exercise}{353}{对称展开与偶次误差}
求 $\lim_{x\to0}[\sqrt{1+x}+\sqrt{1-x}-2]/x^2$。再求实数 $c$，使 $\sqrt{1+x}+\sqrt{1-x}-2-cx^2=o(x^2)$，并计算该差与 $x^4$ 的比值极限。
\end{exercise}
\begin{exercise}{354}{有限极限确定多项式系数}
求使 $\lim_{x\to0}[\ln(1+x)-ax-bx^2]/x^2=2$ 的全部实数 $(a,b)$。对所得参数，继续求分子减去 $2x^2$ 后与 $x^3$ 的比值极限，写出保证运算有效的余项阶数。
\end{exercise}
\begin{exercise}{355}{四阶相消与复合展开}
用泰勒公式求
\[\lim_{x\to0}\frac{1+\tfrac12x^2-\sqrt{1+x^2}}{(\cos x-\mathrm e^{x^2})\sin(x^2)}.\]
再将分母中的 $\mathrm e^{x^2}$ 改为 $\mathrm e^{\alpha x^2}$，讨论极限随实参数 $\alpha$ 的变化，必要时区分左右极限。
\end{exercise}
\begin{exercise}{356}{指数与余弦的四阶比较}
讨论 $\lim_{x\to0}[\cos x-\mathrm e^{x^2/2}]/x^4$，允许结果为无穷大。进一步确定 $a,b$，使 $\cos x-\mathrm e^{ax^2+bx^4}=o(x^4)$，并解释二次项匹配后仍须比较四次项的原因。
\end{exercise}
\begin{exercise}{357}{隐含余项与函数极限恢复}
设 $f$ 在 $0$ 的某去心邻域有定义，且 $\lim_{x\to0}[\sin(6x)+xf(x)]/x^3=0$。不用假设 $f$ 可导，求 $\lim_{x\to0}[6+f(x)]/x^2$。若已知条件中的分母仅为 $x^2$，这一极限是否仍能唯一确定？给出两个具体函数说明。
\end{exercise}
\begin{exercise}{358}{高阶无穷小与参数消去}
求使 $\mathrm e^x-(ax^2+bx+1)=o(x^2)$ 的全部实数 $(a,b)$。对所得参数，求余项的无穷小阶数和非零主项系数；说明 $o(x^2)$ 与 $O(x^2)$ 两种条件对参数的约束有何不同。
\end{exercise}
\begin{exercise}{359}{根式差的阶数与连续消去}
当 $x\to0$ 时，求 $\sqrt{1-2x}-\sqrt[3]{1-3x}$ 相对于 $x$ 的无穷小阶数及主项。进一步求常数 $c$，使减去 $cx^2$ 后成为三阶无穷小，并求相应三阶主项系数。
\end{exercise}
\begin{exercise}{360}{泰勒系数与高阶导数}
设 $f(x)=\ln(1-3x)$。求任意正整数 $n$ 对应的 $f^{(n)}(0)$，并分别用归纳求导和麦克劳林系数验证。写出四阶麦克劳林多项式，说明其在 $0$ 附近的实函数定义范围。
\end{exercise}
\begin{exercise}{361}{泰勒余项的符号}
用带拉格朗日余项的泰勒公式证明 $\mathrm e^x>1+x$ 对 $x\ne0$ 成立，须同时处理 $x<0$。进一步证明对 $x\ge0$ 有 $\mathrm e^x\ge1+x+x^2/2$，并确定等号成立条件。
\end{exercise}
\begin{exercise}{362}{两端泰勒公式与三阶中值}
设 $f\in C^3[-1,1]$，$f(-1)=0$、$f(1)=1$、$f'(0)=0$。分别在 $0$ 处展开至二阶，证明存在 $\xi\in(-1,1)$ 使 $f'''(\xi)=3$。若将 $f'(0)$ 改为 $m$，推导相应三阶导数中值结论，并给出常数 $3$ 的来源。
\end{exercise}
\begin{exercise}{363}{最低点与二阶导数下界}
设 $f\in C^2[0,1]$，$f(0)=f(1)=0$，且在 $[0,1]$ 上的最小值为 $-1$。证明存在 $\xi\in(0,1)$ 使 $f''(\xi)\ge8$。构造满足条件且处处 $f''(x)=8$ 的函数，说明常数 $8$ 无法提高；证明中应明确最低点位于内部的依据。
\end{exercise}
\begin{exercise}{364}{双参数相消与极限分类}
对实数 $a,b$，讨论 $\lim_{x\to0}[\mathrm e^{ax}-1-ax-bx^2]/x^3$ 何时为有限实数，并求其值。再求使分子为 $o(x^3)$ 的全部 $(a,b)$，包括 $a=0$ 的情形。
\end{exercise}
\begin{exercise}{365}{中心差分的误差保证}
设 $f\in C^3[x_0-h,x_0+h]$，$h>0$，且 $|f'''(x)|\le M$。用泰勒公式证明中心差分 $[f(x_0+h)-f(x_0-h)]/(2h)$ 与 $f'(x_0)$ 的误差不超过 $Mh^2/6$。构造一个三次多项式使此界取等号，并讨论仅有 $f''(x_0)$ 存在是否足以保证该误差界。
\end{exercise}
\subsection{函数的单调性与曲线的凹凸性}
\begin{exercise}{366}{分数幂的定义域与单调性}
讨论 $f(x)=x-\tfrac32x^{2/3}$ 在实数范围内的单调区间，指出驻点与不可导点。若将定义域限制为 $[0,+\infty)$，说明单调区间的结论应如何表述，并比较两种定义域下 $x=0$ 附近的行为。
\end{exercise}
\begin{exercise}{367}{逐层求导证明严格不等式}
用导数证明 $\tan x>x+x^3/3$ 对 $0<x<\pi/2$ 成立。进一步证明 $\tan x-x-x^3/3=o(x^3)$，并解释严格正性与高阶小量这两个结论为何可以同时成立。
\end{exercise}
\begin{exercise}{368}{函数比较与商的单调性}
证明当 $x>0$ 时 $\ln(1+x)>\arctan x/(1+x)$。将不等式化为一个在 $0$ 处取零值的函数的正性，再证明该函数在 $(0,+\infty)$ 严格递增，给出导数符号的完整依据。
\end{exercise}
\begin{exercise}{369}{严格单调与导数符号的逻辑关系}
设 $f$ 在 $\mathbb R$ 可导且严格递增。证明 $f'(x)\ge0$，判断能否进一步断言处处 $f'(x)>0$；分别判断 $f(-x)$、$-f(-x)$ 的严格单调性，以及 $f'(-x)$ 的符号。对不成立的加强结论给出具体反例。
\end{exercise}
\begin{exercise}{370}{偶函数方程的根数}
讨论方程 $|x|^{1/4}+|x|^{1/2}-\cos x=0$ 的实根个数。先将根限制在一个有界区间，再利用偶性与正半轴上的单调性证明结论；不以数值近似代替唯一性证明。
\end{exercise}
\begin{exercise}{371}{三次函数与零点个数分岔}
对全部实参数 $a$，讨论方程 $x^3-3x-a=0$ 的不同实根个数，特别求出恰有两个不同实根的参数。给出临界参数下各根的重数，并用单调区间解释根数发生变化的原因。
\end{exercise}
\begin{exercise}{372}{参数约束与高次方程唯一根}
设实数 $a,b,c$ 满足 $3a^2-5b<0$。证明方程 $x^5-2ax^3+3bx+4c=0$ 有且仅有一个实根。若条件放宽为 $3a^2-5b\le0$，判断结论是否仍成立，并处理导数可能为零的情形。
\end{exercise}
\begin{exercise}{373}{严格凸函数的割线斜率}
设 $f\in C^2[a,b]$，$a<b$，且 $f''(x)>0$ 对 $a<x<b$ 成立。令 $F(x)=[f(x)-f(a)]/(x-a)$，$a<x\le b$。证明 $F$ 严格递增，并给出它在 $a$ 处的连续延拓值。若仅知 $f'(x)>0$，该单调性是否仍有保证？给出反例。
\end{exercise}
\begin{exercise}{374}{幂指数比较与参数边界}
设 $a>\mathrm e$，证明对每个 $x>0$ 有 $(a+x)^a<a^{a+x}$。进一步求使这一严格不等式对所有 $x>0$ 成立的全部 $a>0$，并单独处理临界值 $a=\mathrm e$。
\end{exercise}
\begin{exercise}{375}{导数下界与最优统一常数}
设 $\mathrm e<a<b<\mathrm e^2$。证明 $\ln^2b-\ln^2a>4(b-a)/\mathrm e^2$。再证明 $4/\mathrm e^2$ 是使该不等式对全部上述 $a,b$ 成立的最大常数；说明为何开区间中的严格不等号仍允许最优常数由端点极限确定。
\end{exercise}
\begin{exercise}{376}{二阶导数不存在处的凹凸变化}
设 $f(x)=x^{8/3}-x^{5/3}$，其中分数幂按实立方根定义。求定义域、所有凹凸区间和拐点坐标，分别处理二阶导数为零与二阶导数不存在的点。表述凹凸性时同时注明 $f''$ 的符号，避免名称约定造成歧义。
\end{exercise}
\begin{exercise}{377}{绝对值与三角曲线凹凸}
设 $f(x)=|1+\sin x|$。判断它在 $(\pi,3\pi/2)$ 内的严格单调性与凹凸性，写出导数符号依据。再讨论 $x=3\pi/2$ 处的可导性与二阶可导性，说明绝对值能否在全实数范围内直接去掉。
\end{exercise}
\begin{exercise}{378}{严格凸性与中点不等式}
证明对实数 $x\ne y$ 有 $(\mathrm e^x+\mathrm e^y)/2>\mathrm e^{(x+y)/2}$。进一步设 $0<t<1$，用导数或中值定理证明 $t\mathrm e^x+(1-t)\mathrm e^y>\mathrm e^{tx+(1-t)y}$，并完整说明等号出现的条件。
\end{exercise}
\begin{exercise}{379}{分段曲线的凹凸与不可导点}
设 $f(x)=|x^2-3x+2|$。写出分段表达式，求二阶导数存在的全部凹凸区间，并检查 $x=1,2$ 的连续性和可导性。本题将拐点定义为曲线上连续且两侧凹凸性发生改变的点，按此定义求全部拐点，并说明它们是否均满足 $f''(x)=0$。
\end{exercise}
\begin{exercise}{380}{拐点的参数充要条件}
求使 $(0,1)$ 成为曲线 $y=ax^3+bx^2+c$ 拐点的全部实数 $(a,b,c)$。证明条件的充分性与必要性，并解释为何仅有 $f''(0)=0$ 不能排除退化的直线或常函数。
\end{exercise}
\begin{exercise}{381}{偶性与各阶导数的符号}
设 $f\in C^2(\mathbb R)$ 为偶函数，在 $(-\infty,0)$ 有 $f'>0$、$f''<0$。推导 $f'$、$f''$ 的奇偶性，确定它们在 $(0,+\infty)$ 的符号。判断 $0$ 是否为极值点或拐点，并说明依据。
\end{exercise}
\begin{exercise}{382}{三阶判别与水平拐点}
设 $f$ 在 $x_0$ 的邻域三阶可导，$f'(x_0)=f''(x_0)=0$、$f'''(x_0)>0$。证明 $(x_0,f(x_0))$ 为拐点，判断 $x_0$ 是否为极值点。分别通过 $f''$ 的符号变化与泰勒展开说明结论，并指出此处二阶极值判别法不能直接作出判断。
\end{exercise}
\begin{exercise}{383}{微分关系与局部几何形态}
设 $f\in C^2(-\delta,\delta)$，$\delta>0$，满足 $f''(x)+[f'(x)]^2=x$，且 $f'(0)=0$。由该关系证明 $f'''(0)$ 存在并求其值，判断 $0$ 是否为极值点及 $(0,f(0))$ 是否为拐点。论证时须给出足够小邻域内的导数或函数增量符号。
\end{exercise}
\begin{exercise}{384}{函数图形信息与导函数图形}
设 $f(x)=x^3-3x$。先由单调区间和驻点说明 $f'$ 应有的符号与零点，再由凹凸性判断 $f'$ 的单调性，据此手绘 $f'$ 的图形。最后计算 $f'$ 核验，说明 $f$ 的极值点与 $f'$ 的零点之间是否总为一一对应。
\end{exercise}
\begin{exercise}{385}{参数函数的单调区间与严格单调}
设 $f_a(x)=x^3-3ax$，$a\in\mathbb R$。求全部单调区间，确定使 $f_a$ 在整个实数轴上严格递增的参数。对临界参数解释严格递增为何允许导数在孤立点为零，并求该函数的拐点。
\end{exercise}
\begin{exercise}{386}{凸性与零点定位}
设 $f\in C^2[0,1]$，$f(0)=f(1)=0$，且 $f''(x)>0$ 对 $0<x<1$ 成立。证明内部恒有 $f(x)<0$，并证明 $f'$ 在 $(0,1)$ 内恰有一个零点。若将条件改为 $f''\ge0$，判断两个严格结论分别是否仍成立，给出反例或证明。
\end{exercise}
\subsection{函数的极值和最值}
\begin{exercise}{387}{极大值的必要条件与二阶退化}
设 $f$ 在 $x_0$ 的邻域二阶可导，并在 $x_0$ 取得局部极大值。证明 $f'(x_0)=0$、$f''(x_0)\le0$。分别举出二阶导数严格小于零和等于零的例子，说明这两个必要条件能否反过来保证极大值。
\end{exercise}
\begin{exercise}{388}{驻点与不可导点的极值判别}
设 $f(x)=2x-\tfrac32x^{2/3}+1$，分数幂按实立方根定义。求所有驻点、不可导点和极值，用一阶导数符号进行判别。进一步判断该函数在 $\mathbb R$ 上是否有最大值、最小值，说明极值与最值的区别。
\end{exercise}
\begin{exercise}{389}{指数乘积的全局最小值}
求 $f(x)=x2^x$ 的极小值点与极小值，证明该极小值也是全实数范围的最小值。再判断是否存在最大值，并用两端极限支持结论。
\end{exercise}
\begin{exercise}{390}{导函数符号与极值数量}
设 $f$ 在 $\mathbb R$ 可导，且 $f'(x)=x(x+2)(x-1)^2(x-3)$。不用求出 $f$，列出导函数符号表，确定 $f$ 的单调区间以及极大值点、极小值点的数量与位置。说明 $x=1$ 处导数为零却不产生极值的原因，判断仅知导函数能否给出各极值的数值。
\end{exercise}
\begin{exercise}{391}{微分方程约束下的驻点分类}
设 $f$ 在实数轴二阶可导，满足 $xf''(x)+3x[f'(x)]^2=1-\mathrm e^{-x}$。若 $f'(x_0)=0$ 且 $x_0\ne0$，证明 $x_0$ 必为严格极小值点，并排除拐点的可能性。进一步利用关系式在 $0$ 附近的极限讨论 $f''(0)$ 与 $f'(0)$ 满足的等式。
\end{exercise}
\begin{exercise}{392}{导数比值极限与极值}
设 $f$ 在 $a$ 的邻域可导，$f'$ 在 $a$ 连续，且 $\lim_{x\to a}f'(x)/(x-a)=-1$。求 $f'(a)$、$f''(a)$，判断 $a$ 处的极值类型。若已知极限改为 $0$，说明极值类型能否确定，并给出例子。
\end{exercise}
\begin{exercise}{393}{二阶增量极限的信息边界}
设 $f$ 在 $a$ 的邻域有定义，且 $\lim_{x\to a}[f(x)-f(a)]/(x-a)^2=-1$。证明 $f$ 在 $a$ 可导，求 $f'(a)$ 并判断极值类型。判断该极限能否保证 $f''(a)$ 存在；若不能，构造在 $a$ 的邻域可导的反例。
\end{exercise}
\begin{exercise}{394}{分段指数函数与不可导极值点}
设 $f(x)=x^{2x}$（$x>0$）、$f(x)=x\mathrm e^x+1$（$x\le0$）。求各分支导数并检查 $0$ 处的连续性和可导性，确定所有局部极值。再比较各极小值与两端极限，判断全实数范围的最大值和最小值是否存在。
\end{exercise}
\begin{exercise}{395}{隐函数驻点与局部分支}
曲线由 $2y^3-2y^2+2xy-x^2=1$ 给出。在满足隐函数定理条件的点附近把 $y$ 看作 $x$ 的函数，求全部驻点及对应坐标，并判断极值类型。每个候选点均须检验关于 $y$ 的偏导数非零，不能只解 $y'=0$ 的分子。
\end{exercise}
\begin{exercise}{396}{无极值的参数范围}
求使 $f(x)=x^3+3ax^2-ax-1$ 在实数轴上无局部极值的全部实参数 $a$。处理导数有两个实根、重根、无实根三种情况，说明临界参数下驻点为何不构成极值。
\end{exercise}
\begin{exercise}{397}{闭区间最值与根式端点}
求 $f(x)=x+\sqrt{1-x}$ 在 $[-5,1]$ 上的最大值和最小值及全部取到位置。列全内部驻点与端点，并说明 $x=1$ 处导数不存在为何不影响闭区间最值的存在。
\end{exercise}
\begin{exercise}{398}{球外切圆锥的最小体积}
半径 $r>0$ 的球内切于正圆锥，求圆锥高 $h$ 使体积最小，并求最小体积。先由轴截面建立底面半径与 $h,r$ 的关系及 $h$ 的可行范围，再证明所得驻点给出全局最小值，同时检验可行区间两端的体积极限。
\end{exercise}
\begin{exercise}{399}{利润最值与整数产量}
某企业月产量为 $Q$，总成本 $C(Q)=6Q+200$，总收益 $R(Q)=12Q-0.01Q^2$，金额单位均为万元。先把 $Q\ge0$ 视为连续变量，求利润最大时的产量和最大利润，再在 $Q$ 必须为非负整数时核对结论。求企业盈利的产量范围，并区分利润与收益。
\end{exercise}
\begin{exercise}{400}{收入成本与边际量的最优条件}
某企业的总收入为 $R(x)=26x-2x^2-4x^3$，总成本为 $C(x)=8x+x^2$，产量 $x\ge0$ 视为连续变量。求利润函数、边际收入和边际成本，确定利润最大时的产量与最大利润。核验最优点处边际收入等于边际成本，并说明还需哪些符号或边界检查才能证明为全局最大值。
\end{exercise}
\begin{exercise}{401}{共轭指数与等号条件}
设 $p,q>1$ 且 $1/p+1/q=1$。利用函数最值证明对 $x>0$ 有 $x^p/p+1/q\ge x$，确定等号条件。再推出对 $u,v>0$ 有 $uv\le u^p/p+v^q/q$，说明代换过程中使用共轭指数关系的位置。
\end{exercise}
\begin{exercise}{402}{极限给定的切线与严格凸性}
设 $f\in C^2(\mathbb R)$，$\lim_{x\to0}f(x)/x=1$，且 $f''(x)>0$ 对所有实数 $x$ 成立。由极限确定 $f(0)$、$f'(0)$，利用最值或中值定理证明 $f(x)\ge x$，并确定等号条件。若仅有 $f''\ge0$，说明等号条件会怎样变化。
\end{exercise}
\begin{exercise}{403}{约束区间与边界最优}
设 $a\in\mathbb R$，$f_a(x)=x^2-2ax$，生产控制量满足 $0\le x\le1$。求最小值与全部最小值点随 $a$ 的分段表达式，再求最大值和全部最大值点。解释为何只求驻点会遗漏部分参数下的最优解。
\end{exercise}
\begin{exercise}{404}{箱体设计与可行域端点}
从边长为 $L>0$ 的正方形薄板四角各剪去边长为 $x$ 的小正方形，折成无盖长方体。写出体积与严格可行域，求最大体积及对应 $x$，证明全局最优性。若要求底面边长至少为 $L/2$，判断最优方案是否改变并说明依据。
\end{exercise}
\subsection{函数图形的描绘}
\begin{exercise}{405}{水平渐近线与交点}
求 $y=(x+1)/x$ 在 $x\to+\infty$ 与 $x\to-\infty$ 时的水平渐近线，并求铅直渐近线。判断曲线是否与水平渐近线相交，分别说明曲线在两支上从哪一侧逼近该直线。
\end{exercise}
\begin{exercise}{406}{单侧无穷极限与铅直渐近线}
求 $y=1/(1-x^2)$ 的全部铅直渐近线，分别写出每条直线左右两侧的无穷极限。再求水平渐近线，并说明有界的间断点为何不能仅因函数未定义而判为铅直渐近线。
\end{exercise}
\begin{exercise}{407}{约分、可去间断与渐近线}
求 $f(x)=(x^2-x-2)/(x^2-1)$ 的定义域、水平渐近线与铅直渐近线。分别判断 $x=1,-1$ 处的间断类型，并说明代数约分为何不能恢复原函数定义域中缺失的点。
\end{exercise}
\begin{exercise}{408}{斜渐近线与曲线位置}
求 $y=x^2/(x-1)$ 的斜渐近线和铅直渐近线，分别检验正负无穷方向的极限。利用函数与斜渐近线的差，判断两支曲线位于该直线的哪一侧，以及是否存在交点。
\end{exercise}
\begin{exercise}{409}{指数复合函数的单侧渐近线}
求 $y=(2x-1)\mathrm e^{1/x}$ 的全部水平、铅直和斜渐近线。重点分别计算 $x\to0^+$、$x\to0^-$ 的极限与 $x\to\pm\infty$ 时的斜截距，说明同一铅直渐近线的两侧是否都必须发散。
\end{exercise}
\begin{exercise}{410}{振荡因子与渐近线判定}
对定义在 $x>0$ 的函数 $f(x)=x\sin(1/x)$，求全部水平与铅直渐近线，并讨论 $x\to0^+$ 时能否连续补值。说明 $\sin(1/x)$ 的振荡为何不妨碍 $f$ 在右端点有有限极限。
\end{exercise}
\begin{exercise}{411}{两端斜渐近线与单侧发散}
求 $f(x)=x\mathrm e^{1/x}$ 的全部水平、铅直和斜渐近线。写出 $0$ 两侧的极限，并证明在正负无穷处得到的斜渐近线相同。
\end{exercise}
\begin{exercise}{412}{对数指数函数的渐近线条数}
求 $f(x)=\ln(1+\mathrm e^x)/x$ 的全部渐近线，分别处理 $x\to0^\pm$、$x\to\pm\infty$。对正无穷方向先抽出对数中的主导指数，再说明是否还可能存在斜率非零的斜渐近线。
\end{exercise}
\begin{exercise}{413}{反正切复合与有限间断辨析}
求 $f(x)=\mathrm e^{1/x}\arctan[(x^2-x+2)/((x-1)(x+2))]$ 的定义域及全部渐近线，给出渐近线总条数。分别检查 $x=0,1,-2$ 两侧的极限，区分无穷间断和有限跳跃，并计算正负无穷方向的极限。
\end{exercise}
\begin{exercise}{414}{有理函数的完整图形分析}
对 $f(x)=(x-1)^3/(x+1)^2$，求定义域、截距、全部渐近线、单调区间、极值、凹凸区间和拐点，注明关键点及两端趋势，据此手绘图形。用精确计算说明曲线是否能够穿过其斜渐近线；若能，求交点。
\end{exercise}
\begin{exercise}{415}{周期函数的图形描绘}
对 $f(x)=\cos x/\cos2x$，确定定义域、奇偶性与最小正周期。在 $[-\pi,\pi]$ 内求全部铅直渐近线、与坐标轴的交点、单调区间、极值、凹凸区间和拐点，据此手绘该区间图形并说明如何延拓至全定义域。
\end{exercise}
\begin{exercise}{416}{渐近线可被穿越}
设 $f(x)=x+\sin x/x$，$x\ne0$。求它的全部渐近线，并判断 $0$ 处的间断类型。求与斜渐近线的全部交点，证明在任意充分大的正半轴尾部仍存在交点，说明渐近线的定义实际要求的极限是什么。
\end{exercise}
\begin{exercise}{417}{参数抵消与渐近线变化}
设 $f_a(x)=(x^2+ax+1)/(x-1)$，定义域为 $\mathbb R\setminus\{1\}$。求全部渐近线随 $a$ 的变化，并确定何时 $x=1$ 为可去间断点。对非退化参数，判断曲线能否与其斜渐近线相交，证明结论。
\end{exercise}
\subsection{曲率}
\begin{exercise}{418}{显式与参数表示的曲率一致性}
求曲线 $xy=1$ 在 $(1,1)$ 处的曲率与曲率半径，分别用显式表示 $y=1/x$ 和参数表示 $(x,y)=(\mathrm e^t,\mathrm e^{-t})$ 验证。说明曲率为何与参数行进的快慢无关。
\end{exercise}
\begin{exercise}{419}{星形线的正则参数与曲率半径}
设 $a>0$，曲线为 $x=a\cos^3t$、$y=a\sin^3t$。求满足 $\sin t_0\cos t_0\ne0$ 的点处的曲率和曲率半径，确定第一象限弧上曲率的最小值。讨论 $t_0=k\pi/2$ 时通常的正则参数曲率公式为何不能直接使用，并求沿相邻弧趋近时的曲率极限。
\end{exercise}
\begin{exercise}{420}{法向方向与曲率圆}
求抛物线 $y=2x^2-3x+3$ 在 $M(1,2)$ 处的曲率、曲率中心与曲率圆方程。证明圆心位于该点的法线上，并核验该圆经过 $M$；说明仅知道曲率半径为何还不能确定圆心的位置。
\end{exercise}
\begin{exercise}{421}{弧长变量的链式求导}
设 $\rho(x)$ 是抛物线 $y=\sqrt x$ 在 $M(x,\sqrt x)$、$x\ge1$ 处的曲率半径，$s(x)$ 是从 $A(1,1)$ 到 $M$ 沿曲线的弧长。求 $\rho(x)$、$\dd s/\dd x$，再计算 $3\rho\,\dd^2\rho/\dd s^2-(\dd\rho/\dd s)^2$。明确二阶弧长导数的链式法则，不能把它直接替换为关于 $x$ 的二阶导数。
\end{exercise}
\begin{exercise}{422}{抛物线最大曲率与长度缩放}
设 $a>0$，曲线为 $y=ax^2$。求曲率函数、最大曲率及对应点，写出该点的曲率圆。若把曲线上所有坐标同时放大 $\lambda>0$ 倍，证明对应点的曲率变为原来的 $1/\lambda$，并解释量纲关系。
\end{exercise}
\begin{exercise}{423}{椭圆曲率的极值}
设 $a>b>0$，椭圆由 $x=a\cos t$、$y=b\sin t$ 表示。求曲率函数及其最大、最小值，给出对应点和曲率半径。再讨论 $a=b$ 时结论的变化，并将结果与圆的几何曲率核对。
\end{exercise}
